% Generated draft: requires mathematical and source review.
\documentclass[11pt]{article}
\usepackage[margin=25mm]{geometry}
\usepackage{fontspec}
\setmainfont{Latin Modern Roman}
\usepackage{amsmath,amssymb,mathtools,unicode-math}
\setmathfont{Latin Modern Math}
\usepackage{enumitem,needspace,xurl,hyperref}
\setlength{\parindent}{0pt}
\setlength{\parskip}{5pt}
\setlength{\emergencystretch}{3em}
\newcommand{\N}{\mathbb{N}}
\newcommand{\Z}{\mathbb{Z}}
\newcommand{\Q}{\mathbb{Q}}
\newcommand{\R}{\mathbb{R}}
\newcommand{\C}{\mathbb{C}}
\newcommand{\F}{\mathbb{F}}
\newcommand{\A}{\mathbb{A}}
\newcommand{\PP}{\mathbb{P}}
\newcommand{\E}{\mathbb{E}}
\newcommand{\eps}{\varepsilon}
\newcommand{\dd}{\,\mathrm{d}}
\begin{document}
\section*{UnsolvedMath ID: 4000009}
\textbf{Source problem code:} catalogue entry
\subsection{Definitions and mathematical statement}
Let $(X,d)$ be a metric space with Markov transition measures $m_x$ and an invariant probability measure $\nu$. For transition measures with finite first moments, write
\[
\kappa(x,y)=1-\frac{W_1(m_x,m_y)}{d(x,y)}\quad(x\ne y),
\]
where $W_1$ is the infimum of the expected distance over couplings of its arguments. Positive coarse Ricci curvature means $\kappa(x,y)\geq\kappa_0>0$ for all distinct $x,y$ in the relevant space. A usual local spread convention, when finite, is
\[
\sigma_x^2=\frac12\int_X\int_X d(u,v)^2\,dm_x(u)\,dm_x(v).
\]
The exact local-variance normalization and boundedness hypothesis intended in the source are not supplied.

Find weaker local assumptions than the bounded-local-variance hypothesis in the positive-curvature concentration theorem, while obtaining concentration estimates governed by the actual function $x\mapsto\sigma_x^2$. In an invariant-measure formulation, the desired estimates control tails of integrable $1$-Lipschitz observables $F$, such as
\[
\nu\{x:F(x)-\nu(F)\geq t\},\qquad t\geq0,
\]
using curvature and local variance rather than a coarse worst-case substitute. The bounds should also remain informative with finite, nondiverging constants when a discrete transition kernel approaches a continuous-time process and individual transition probabilities become small. A scale family, assumptions on its limiting generator, and the precise tail function are not given, so this is a request for a theorem with suitable hypotheses rather than a specified numerical inequality.

The supplied summary reports local refinements in some settings, failure of naive replacements by local spread for small transition probabilities (including binomial measures on the cube), and deterioration in a continuous-time limit. It does not provide the full existing concentration theorem or an exact proposed replacement.
\subsection{Short English statement}
Under positive coarse Ricci curvature, seek concentration estimates that use actual local variance under weaker local assumptions. The estimates should retain meaningful finite constants as discrete transitions converge to continuous time.
\subsection{Sources}
\begin{itemize}
\item Yann Ollivier, \emph{Discrete Ricci curvature: Open problems} (2008), Problem I, \url{https://www.yann-ollivier.org/rech/publs/problems_curvmarkov.pdf}; accessed 2026-07-29, with PDF-text extraction recorded in the supplied provenance.
\item The supplied research summary refers to Ollivier's main article, Gromov--Milman-style concentration, and local-variance refinements without a title or theorem citation. It schematically writes $\exp(-t\sqrt{\kappa}\,\sigma_\infty)$; the definitions, normalizations, hypotheses, and exact placement of the scale parameter are not supplied, so that expression is not adopted here as a precise theorem.
\item Local registry \texttt{data/unsolvedmath/problems.json}, permanent UnsolvedMath ID $4000009$, source code catalogue entry, \emph{Local assumptions for concentration}. Dataset attribution: Ulam AI and the UnsolvedMath Contributors.
\end{itemize}
\end{document}
